\chapter{Limitations and Future Work}
\label{ch:limitations}

\section{Limitations}

\subsection{Synthetic Data}

All results in this report are derived from synthetic simulation data. While the simulation models are based on realistic physical equations and empirical load profiles, they cannot capture the full complexity of real-world conditions: irregular cloud patterns, occupant behavior variability, appliance-level consumption signatures, and sensor noise characteristics. The PV forecast results ($R^2 \approx 1.0$) are a direct consequence of using deterministic synthetic data and should not be interpreted as expected real-world performance.

\subsection{Single Building}

The system is designed for a single small building. Scaling to multiple buildings, a campus, or a microgrid would require additional components: inter-building load balancing, shared storage management, and aggregated optimization.

\subsection{No Grid Export}

The current model assumes grid import only ($P_{\text{grid}} \geq 0$). Morocco's net metering regulations allow feed-in for installations up to 5~MW, but the economic value of exported energy depends on the specific contract and was not modeled.

\subsection{Simplified Tariff}

The economic analysis uses a flat tariff estimate (1.20~MAD/kWh) with a peak multiplier. The actual ONEE residential tariff uses consumption tranches (progressive pricing), which would change the cost calculations. The tariff values should be verified from official ONEE publications before deployment.

\subsection{Anomaly Detection Precision}

The statistical anomaly detector achieves high recall (71\%) but low precision (8.5\%) on synthetic data, generating many false positives. With real data and tuned parameters ($k > 3$, contextual filtering), precision is expected to improve.

\subsection{Battery Degradation Model}

The battery degradation cost is modeled as a simple linear function of discharged energy. A more accurate model would account for depth of discharge, C-rate, temperature, and calendar aging effects.

\section{Future Work}

\begin{enumerate}
    \item \textbf{Real sensor data:} Deploy the ESP32 prototype at an actual building and retrain ML models on real measurements.

    \item \textbf{Grid export / net metering:} Extend the optimization model to allow bidirectional grid flow and model feed-in compensation.

    \item \textbf{Weather API integration:} Use forecasted weather data (e.g., OpenWeatherMap) instead of historical irradiance for day-ahead PV prediction.

    \item \textbf{Electric vehicle integration:} Model an EV as both a flexible load and a mobile storage resource (vehicle-to-home).

    \item \textbf{Reinforcement learning:} Once sufficient real operational data is available, train an RL agent as a third optimization strategy and compare with MILP.

    \item \textbf{Multi-building aggregation:} Extend to campus-level energy management with shared storage and peer-to-peer energy trading.

    \item \textbf{Edge ML on ESP32:} Deploy a lightweight anomaly detection model directly on the microcontroller for local, low-latency alerts.

    \item \textbf{Mobile application:} Develop a mobile app (Flutter or React Native) for operator notifications and remote monitoring.

    \item \textbf{Full-year optimization:} Run the MILP optimizer over the complete 365-day dataset to provide more accurate annual performance projections.
\end{enumerate}
